\chapter{Zur Passionszeit}
\quellseite{1}

„Gott aber, was er durch den Mund aller seiner Propheten verkündiget hat, wie Christus leiden sollte, hat’s also erfüllt.“ Apost. 3, 18.

Die heilige Passionszeit wird von der Lutherischen Kirche zu dem Zwecke gefeiert, daß in ihr mehr als sonst im Jahre das bittere Leiden unsers HErrn JEsu Christi betrachtet werde. Hierbei ist es wichtig, auch gerade die Weissagungen der heiligen Propheten über das Leiden des zukünftigen Messias zu erwägen. Nur so kann man gewiß werden, daß JEsus Christus derjenige ist, über den Gott zu Adam, zu den Erzvätern, zum Volke Gottes im alten Testament geredet hat, daß man die Worte der Propheten von dem zukünftigen Heiland mit der Lebensgeschichte Christi vergleicht und zusihet, ob auch alles genau in Erfüllung gegangen ist. Diese Weise, zur Gewißheit des Glaubens zu kommen, lauschen wir unserm HErrn JEsu selbst ab. Als er nämlich von den Todten auferstanden war, strafte er die verstörten, nach Emmaus gehenden Jünger, daß sie so träge zu glauben gewesen waren, und sprach: „Mußte nicht Christus solches leiden und zu seiner Herrlichkeit eingehen? Und fing an von Mose und allen Propheten, und legte ihnen alle Schriften aus, die von ihm gesagt waren“. Luc. 24, 25—27. Und als er später im Kreise seiner Jünger stand, schlug er eben denselben Weg ein, sie zum fröhlichen Glauben an ihn als an den verheißenen Heiland zu führen. Er sprach zu ihnen: „Das sind die Reden, die ich zu euch sagte, da ich noch bei euch war; denn es muß alles erfüllet werden, was von mir geschrieben ist im Gesetz Mosis, in den Propheten, und in den Psalmen. Da öffnete er ihnen das Verständniß, daß sie die Schrift verstunden“, Luc. 24, 44. 45. Als dann die Apostel mit der Predigt des Evangeliums in alle Welt auszogen, kamen auch sie in ihren Predigten beständig auf das Wort der Propheten vom zukünftigen Messias zurück und zeigten: Das ist alles bis ins Kleinste an JEsu von Nazareth erfüllt.

So mußte denn ein jeder, der aufrichtig forschte und sich nicht muthwillig gegen das Zeugniß der heiligen Schrift und der heiligen Apostel verschloß, erkennen: JEsus ist der Messias, der verheißene Christ Gottes.

Im Folgenden wollen wir uns den Weissagungen über das Leiden und Sterben des Messias zuwenden und achten, wie genau sie in JEsu von Nazareth erfüllt sind.

\gesperrt{Der Messias soll auf einem Eselsfüllen in Jerusalem einziehen.} Sacharja ruft: „Du Tochter Zion, freue dich sehr, und du Tochter Jerusalem, jauchze; siehe, dein König kommt zu dir, ein Gerechter und ein Helfer, arm, und reitet auf einem Esel, und auf einem jungen Füllen der Eselin“, Sach. 9, 9. — Das ist an Christo erfüllt. Matthäus berichtet: „Und die Jünger brachten die Eselin und das Füllen und legten ihre Kleider darauf und setzten ihn darauf.... Und als er zu Jerusalem einzog, erregte sich die ganze Stadt und sprach: Wer ist der? Das Volk aber sprach: Das ist der JEsus, der Prophet von Nazareth aus Galiläa“, Matth. 21, 7. ff.

\gesperrt{Die Obersten werden rathschlagen, den Messias zu tödten.} Im 2. Psalm heißt es: „Warum toben die Heiden, und die Leute reden so vergeblich? Die Könige im Lande lehnen sich auf, und die Herren rathschlagen mit einander wider den HErrn und seinen Gesalbten“, Ps. 2, 1. 2. — Das ist an Christo erfüllt. Johannes erzählt: „Da versammelten die Hohenpriester und die Pharisäer einen Rath und sprachen: Was thun wir? Dieser Mensch thut viel Zeichen. Einer aber unter ihnen, Caiphas, der desselben Jahrs Hoherpriester war, sprach zu ihnen: Es ist uns besser, Ein Mensch sterbe für das Volk, denn daß das ganze Volk verderbe. Von dem Tage an rathschlagten sie, wie sie ihn tödteten“, Joh. 11, 47. 49. 53. Apost. 4, 23—27.

\gesperrt{Der Messias wird durch einen Jünger an die Feinde verrathen werden.} Im 41. Psalm klagt er: „Sie haben ein Bubentstück über mich beschlossen: Wenn er liegt, soll er nicht wieder aufstehen. Auch mein Freund, dem ich
\quellseite{2}
mich vertraute, der mein Brod aß, tritt mich unter die Füße“, B. 9, 10. — Das ist an Christo erfüllt. Matthäus sagt: „Und da JEsus noch redete, siehe, da kam Judas, der Zwölfen einer, und mit ihm eine große Schaar, mit Schwertern und mit Stangen, von den Hohenpriestern und Aeltesten des Volks. JEsus aber sprach zu ihm: Mein Freund, warum bist du kommen? Juda, verräthest du des Menschen Sohn mit einem Kuß?“ Matth. 26, 47. ff. Luc. 22, 48. Joh. 13, 18.

\gesperrt{Der gefangene Messias wird von seinen Jüngern verlassen werden.} Sacharja ruft: „Schwert, mache dich auf über meinen Hirten und über den Mann, der mir der nächste ist, spricht der HErr Zebaoth. Schlage den Hirten, so wird die Heerde sich zerstreuen“, Sach. 13, 7. — Das ist an Christo erfüllt. Marcus erzählt: „Die (Feinde) aber legten ihre Hände an ihn und griffen ihn. Und die Jünger verließen ihn alle und flohen“, Marc. 14, 46. 50. Vgl. B. 27.

\gesperrt{Der Messias wird um dreißig Silberlinge verkauft werden.} Sacharja weissagt: „Gefällt es euch, so bringet her, wie viel ich gelte; wo nicht, so laßt es anstehen. Und sie wogen dar, wie viel ich galt, dreißig Silberlinge“, Sach. 11, 12. — Das ist an Christo erfüllt. Judas sprach zu den Hohenpriestern: „Was wollt ihr mir geben? Ich will ihn euch verrathen. Und sie boten ihm dreißig Silberlinge“, Matth. 26, 15.

\gesperrt{Eines Töpfers Acker wird um die zurückgebrachten Silberlinge gekauft werden.} Sacharja verkündigt: „Und der HErr sprach zu mir: Wirf’s hin, daß es dem Töpfer gegeben werde. Ei, eine treffliche Summa, der ich werth geachtet bin von ihnen! Und ich nahm die dreißig Silberlinge, und warf sie ins Haus des HErrn, daß dem Töpfer gegeben würde“, Sach. 11, 13. — Das ist erfüllt. Matthäus berichtet: „Da das sahe Judas, der ihn verrathen hatte, daß er verdammt war zum Tode, gereuete es ihn, und brachte herwieder die dreißig Silberlinge den Hohenpriestern und den Aeltesten. Und er warf die Silberlinge in den Tempel. Sie hielten aber einen Rath und kauften einen Töpfersacker darum, zum Begräbniß der Pilger“, Matth. 27, 3—9.

\gesperrt{Der Verräther des Messias wird eines jähen Todes sterben.} Im 109. Psalm wird gesagt: „Seiner Tage müssen wenig werden, und sein Amt müsse ein anderer empfahen“, B. 8. — Das ist am Verräther Christi erfüllt. Beim Evangelisten Matthäus heißt es: „Und Judas warf die Silberlinge in den Tempel, hub sich davon, ging hin und erhengete sich selbst“, Cap. 27, 5. Apost. 1, 20.

\gesperrt{Der Messias wird auf das schmählichste mißhandelt werden.} Jesaias beschreibt es so: „Ich hielt meinen Rücken dar denen, die mich schlugen, und meine Wangen denen, die mich rauften; mein Angesicht verbarg ich nicht vor Schmach und Speichel. Ich habe mein Angesicht dargeboten als einen Kieselstein“, Jes. 50, 6. 7. — Das ist an Christo erfüllt. Marcus erzählt: „Und Pilatus überantwortete ihnen JEsum, daß er gegeißelt und gekreuziget würde. Die Kriegsknechte aber führeten ihn hinein in das Richthaus und riefen zusammen die ganze Schaar und zogen ihm einen Purpur an und flochten eine dornene Krone und setzten sie ihm auf; und fingen an, ihn zu grüßen: Gegrüßet seist du, der Juden König. Und schlugen ihn mit dem Rohr und verspieen ihn, und fielen auf die Kniee und beteten ihn an“, Marc. 15, 15—19. Matth. 26, 67.

\gesperrt{Der Messias wird die Mißhandlung still erdulden.} Jesaias sagt: „Da er gestraft und gemartert ward, that er seinen Mund nicht auf, wie ein Lamm, das zur Schlachtbank geführet wird, und wie ein Schaf, das verstummet vor seinem Scherer, und seinen Mund nicht aufthut“, Jes. 53, 7. — Das ist an Christo erfüllt. Matthäus bezeugt: „Und da er verklagt ward von den Hohenpriestern und Aeltesten, antwortete er nichts. Da sprach Pilatus: Hörest du nicht, wie hart sie dich verklagen? Und er antwortete ihm nicht auf ein Wort, also, daß sich auch der Landpfleger sehr verwunderte“, Matth. 27, 12. ff.

\gesperrt{Der Messias soll unter Uebelthätern gekreuziget werden.} Der 22. Psalm weissagt: „Sie haben meine Hände und Füße durchgraben. Alle meine Gebeine haben sich zertrennet“, Ps. 22, 15. 17. Und Jesaias sagt: „Er ist unter die Uebelthäter gerechnet“, Jes. 53, 12. — Das ist an Christo erfüllt. Lucas erzählt: „Es wurden aber auch hingeführt zween andere Uebelthäter, daß sie mit ihm abgethan würden. Und als sie kamen an die Stätte, die da heißt Schädelstätte, kreuzigten sie ihn daselbst, und die Uebelthäter mit ihm, einen zur Rechten und einen zur Linken“, Luc. 23, 32. 33. Marc. 15, 27. 28.

\gesperrt{Die Kleider des Messias sollen vertheilt und verlost werden.} Der 22. Psalm weissagt: „Sie theilen meine Kleider unter sich und werfen das Loos um mein Gewand“, Ps. 22, 19. — Das ist an Christo erfüllt. Lucas meldet: „Die Kriegsknechte aber, da sie JEsum gekreuziget hatten, nahmen sie seine Kleider, und machten vier Theile, einem jeglichen Kriegsknecht ein Theil, dazu auch den Rock. Der Rock aber war ungenähet, von oben an gewirkt durch und durch. Da sprachen sie unter einander: Lasset uns den nicht zertheilen, sondern darum losen, weß er sein soll. Solches thaten die Kriegsknechte“, Luc. 19, 23. 24.

\gesperrt{Der gekreuzigte Messias wird für seine Peiniger beten.} Bei Jesaias heißt es: „Er hat für die Uebelthäter gebeten“, Jes. 53, 12. — Das ist in Christo erfüllt. Lucas berichtet: „JEsus aber sprach: Vater, vergib ihnen; denn sie wissen nicht, was sie thun“, Luc. 23, 34.

\gesperrt{Der gekreuzigte Messias soll verhöhnt werden.} Im 22. Psalm klagt er: „Ich bin ein Wurm und kein Mensch, ein Spott der Leute und Verachtung des Volks. Alle, die mich sehen, spotten mein, sperren das Maul auf und schütteln den Kopf: Er klage es dem HErrn, der helfe ihm aus und errette ihn, hat er Lust zu ihm“, Ps. 22, 7—9. — Das ist an Christo erfüllt. Matthäus bezeugt: „Die aber vorübergingen, lästerten ihn und schüttelten ihre Köpfe, und sprachen: Der du den Tempel Gottes zerbrichst, und bauest ihn in dreien Tagen, hilf dir selber. Bist du Gottes Sohn, so steig herab vom Kreuz. Er hat Gott vertrauet, der erlöse ihn nun, lüstet’s ihn“, Matth. 27, 39. 40. 43.

\gesperrt{Der Messias wird von Gott verlassen sein.} Im 22. Psalm ruft er aus: „Mein Gott, mein Gott, warum hast du mich verlassen? Ich heule, aber meine Hülfe ist ferne“, Ps. 22, 2. — Das ist in Christo erfüllt. Matthäus erzählt: „Und
\quellseite{3}
in die neunte Stunde schrie JEsus laut und sprach: Eli, Eli, lama asabthani? das ist: Mein Gott, mein Gott, warum hast du mich verlassen? Matth. 27, 46.

\gesperrt{Der Messias wird in seinem Durst mit Galle und Essig getränket werden.} In den Psalmen klagt er: „Meine Kräfte sind vertrocknet wie ein Scherben, und meine Zunge klebet an meinem Gaumen“, Ps. 22, 16. „Und sie geben mir Galle zu essen und Essig zu trinken in meinem großen Durst“, Ps. 69, 22. — Das ist an Christo erfüllt. Der Evangelist sagt: „Sie gaben JEsu Essig zu trinken mit Galle vermischt“, Matth. 27, 34. „Darnach, als JEsus wußte, daß schon alles vollbracht war, daß die Schrift erfüllet würde, spricht er: Mich dürstet. Da stund ein Gefäß voll Essigs. Sie aber fülleten einen Schwamm mit Essig und legten ihn um einen Ysop und hielten es ihm dar zum Munde“, Joh. 19, 28. 29.

\gesperrt{Der Messias wird sterben, aber freiwillig.} Jesaias verkündigt: „Er ist aus dem Lande der Lebendigen weggerissen.“ „Er hat sein Leben in den Tod gegeben“, Jes. 53, 8. 12. — Das ist in Christo erfüllt. Christus spricht: „Niemand nimmt mein Leben von mir, sondern ich lasse es von mir selber. Ich habe es Macht zu lassen“, Joh. 10, 18. Darum berichtet Lucas: „Und JEsus rief laut und sprach: Vater, ich befehle meinen Geist in deine Hände. Und als er das gesagt, verschied er“, Luc. 23, 46.

\gesperrt{Die Seite des gekreuzigten Messias soll durchstochen werden.} Sacharja weissagt: „Sie werden mich ansehen, welchen sie zerstochen haben“, Sach. 12, 10. — Das ist an Christo erfüllt. Johannes bezeuget: „Als die Kriegsknechte zu JEsu kamen, da sie sahen, daß er schon gestorben war, brachen sie ihm die Beine nicht; sondern der Kriegsknechte einer öffnete seine Seite mit einem Speer, und alsbald ging Blut und Wasser heraus“, Joh. 19, 33. 34.

\gesperrt{Der Messias soll ehrenhaft bestattet werden.} Der 22. Psalm sagt: „Und du legest mich in des Todes Staub“, Ps. 22, 16., und Jesaias spricht: „Und er ist begraben wie die Gottlosen und gestorben wie ein Reicher“ (das heißt, die Juden gedachten ihm sein Grab zu geben bei den Gottlosen, aber im Tode war er bei einem Reichen), Jes. 53, 9. — Das ist in Christo erfüllt. Matthäus erzählt: „Am Abend aber kam ein reicher Mann von Arimathia, der hieß Joseph, welcher auch ein Jünger JEsu war. Der ging zu Pilato und bat um den Leib JEsu. Da befahl Pilatus, man sollte ihm ihn geben. Und Joseph nahm den Leib und wickelte ihn in eine reine Leinwand, und legte ihn in sein eigen neu Grab“, Matth. 27, 57—60.

\gesperrt{Der Messias wird unschuldig gekreuzigt und getödtet werden.} Jesaias sagt: „Er hat niemand Unrecht gethan, noch ist Betrug in seinem Munde gewesen“, Jes. 53, 9. — Das ist in Christo erfüllt. „Pilatus sprach zu ihnen: Ihr habt diesen Menschen zu mir gebracht, als der das Volk abwendig. Und siehe, ich hab ihn vor euch gehöret, und finde an dem Menschen der Sachen keine, der ihr ihn beschuldiget; Herodes auch nicht; denn ich habe euch zu ihm gesandt, und siehe, man hat nichts auf ihn gebracht, das des Todes werth sei. Darum will ich ihn züchtigen und loslassen“, Luc. 23, 13—16.

\gesperrt{Des Messias Leiden soll sein zur Erlösung der Welt.} Jesaias sagt: „Er ist um unserer Missethat willen verwundet und um unserer Sünde willen zerschlagen. Die Strafe liegt auf ihm, auf daß wir Friede hätten, und durch seine Wunden sind wir geheilet“, Jes. 53, 5. — Das ist in Christo erfüllt. Der Apostel bezeuget: „Christus hat uns erlöset von dem Fluch des Gesetzes, da er ward ein Fluch für uns (denn es stehet geschrieben: Verflucht sei jedermann, der am Holz hänget)“, Gal. 3, 13. „Christus hat sich selbst gegeben für alle zur Erlösung“, 1 Tim. 2, 6.

Wie wunderbar genau hat also der Heilige Geist die Leiden des Messias vorhergesagt, und wie wunderbar getreu ist jede Weissagung in Christo JEsu erfüllt! Gott hat aber absichtlich auch die kleinsten Nebenumstände aufs deutlichste vorherverkündigt, weil unter allen Millionen Menschen nur Ein Mittler zwischen Gott und den Menschen, nur Ein Erlöser sein konnte und sollte. Gott wollte so untrügliche Merkmale dieses Heilandes angeben, damit, wenn er nun wirklich ihn senden würde, alle Welt erkennen müßte: Ja, der ist’s, von dem Gott durch die Propheten geredet hat; der ist der Helfer aus Sünde und Tod! Nun hat kein Mensch ein Recht, an Christo zu zweifeln, ob er auch wirklich der zukünftige Erlöser ist, denn siehe doch, in ihm ist alle Schrift erfüllt.

\begin{verse}
O Lamm Gottes unschuldig\\
Am Stamm des Kreuzes geschlachtet,\\
Allzeit funden geduldig,\\
Wiewohl du wärest verachtet,\\
All Sünd hast du getragen,\\
Sonst müßten wir verzagen.\\
Erbarm dich unser, o JEsu!
\end{verse}

Friedr. Brand.
